% Atur variabel berikut sesuai dengan identitas Anda

% nama
\newcommand{\name}{Hernanda Achmad Priyatna}
\newcommand{\authorname}{Priyatna, Hernanda Achmad}
\newcommand{\nickname}{Hernanda}
\newcommand{\advisor}{Dr. Rudy Dikairono, S.T., M.T.}
\newcommand{\coadvisor}{Muhammad Attamimi, B.Eng., M.Eng., Ph.D.}
\newcommand{\examinerone}{}
\newcommand{\examinertwo}{}
\newcommand{\examinerthree}{}
\newcommand{\headofdepartment}{}

% identitas
\newcommand{\nrp}{5022221114}
\newcommand{\advisornip}{198103252005011002}
\newcommand{\coadvisornip}{198503272019031006}
\newcommand{\examineronenip}{000000000000000000}
\newcommand{\examinertwonip}{000000000000000000}
\newcommand{\examinerthreenip}{000000000000000000}
\newcommand{\headofdepartmentnip}{000000000000000000}

% judul
\newcommand{\tatitle}{NAVIGASI MOBIL OTONOM MENGGUNAKAN SEGMENTASI VISUAL DAN DATA GPS UNTUK MANUVER DI JALAN RAYA ITS}
\newcommand{\engtatitle}{\emph{NAVIGATION OF AN AUTONOMOUS VEHICLE USING VISUAL SEGMENTATION AND GPS DATA FOR MANEUVERING ON ITS ROADWAYS}}

% tempat
\newcommand{\place}{Surabaya}

% jurusan
\newcommand{\studyprogram}{S1 Teknik Elektro}
\newcommand{\engstudyprogram}{Undergraduate Study Program of Electrical Engineering}

% fakultas
\newcommand{\faculty}{Fakultas Teknologi Elektro dan Informatika Cerdas}
\newcommand{\engfaculty}{Faculty of Electrical Technology and Intelligent Information Technology}

% singkatan fakultas
\newcommand{\facultyshort}{FTEIC}
\newcommand{\engfacultyshort}{FTEIC}

% departemen
\newcommand{\department}{Departemen Teknik Elektro}
\newcommand{\engdepartment}{Department of Electrical Engineering}

% kode mata kuliah
\newcommand{\coursecode}{EE234799}
